\documentclass[11pt,a4paper]{article}
\usepackage[utf8]{inputenc}
\usepackage[T1]{fontenc}
\usepackage{amsmath,amssymb,amsfonts}
\usepackage{booktabs}
\usepackage{multirow}
\usepackage{graphicx}
\usepackage{xcolor}
\usepackage[margin=2.5cm]{geometry}
\usepackage{hyperref}
\usepackage{enumitem}
\usepackage{caption}
\usepackage{subcaption}

\title{Per-Head KV Cache Quantization:\\When It Helps, When It Hurts, and How to Tell}

\author{Anonymous}

\date{}

\begin{document}

\maketitle

%% ========== ABSTRACT ==========
\begin{abstract}
Per-head mixed-precision KV cache quantization promises better quality--compression tradeoffs than uniform bit-width allocation.
We conduct a systematic empirical study across five diverse architectures (Qwen2.5-7B, Llama-3.1-8B, Mistral-7B-v0.3, Yi-1.5-9B, Gemma-2-9B) spanning GQA ratios from 2:1 to 8:1, evaluating per-head allocation on hardware-aligned bit-widths \{2,4,8,16\}.
Our central finding is that \textbf{non-uniform allocation benefits are rare}: at 4$\times$ compression, only 1 of 5 architectures (Qwen2.5-7B) shows meaningful improvement over uniform 4-bit (+29.4\% perplexity reduction), while the remaining four show no benefit or minor degradation ($-$0.4\% to +0.2\%).
We provide three resource contributions:
(1)~a \textbf{systematic benchmark} quantifying per-head allocation outcomes across architectures, compression ratios, and evaluation settings;
(2)~a \textbf{sensitivity-based taxonomy} classifying architectures into Spread, Spike, and Flat categories that predict allocation outcomes;
and (3)~a \textbf{two-metric diagnostic} (top-5\% concentration $\gtrsim$0.50 AND $t_1/t_5 \lesssim 0.70$) that correctly classifies all tested models, enabling practitioners to decide whether profiling is worthwhile before investing GPU time.
We additionally characterize a compounding failure mode where unconstrained solvers over-allocate 2-bit heads, and show that a simple cardinality constraint prevents it---reducing the problem to where solver choice is irrelevant.
We report honest negative results, including downstream task degradation on 4/5 architectures, 8\% inference latency overhead from mixed-precision, and profiling scalability limitations at 70B+ scale.
\end{abstract}

%% ========== INTRODUCTION ==========
\section{Introduction}

KV cache quantization is a widely adopted technique for reducing memory consumption during autoregressive inference in large language models~\cite{hooper2024kvquant,liu2024kivi}.
Rather than storing key--value activations in full precision, practitioners quantize them to lower bit-widths, trading quality for memory.
A natural extension is \emph{per-head mixed-precision} allocation: assigning different bit-widths to different attention heads based on their quantization sensitivity~\cite{zhang2024ratequant}.
The implicit assumption is that non-uniform allocation always improves upon uniform quantization---heads with low sensitivity can tolerate aggressive compression while sensitive heads retain higher precision.

In this work, we test this assumption systematically.
We evaluate per-head allocation on hardware-aligned bit-widths \{2,4,8,16\} across five architectures spanning diverse GQA configurations (2:1 to 8:1), multiple compression ratios (2$\times$, 4$\times$, 8$\times$), and both perplexity and downstream task metrics.
Our findings challenge the prevailing assumption:

\begin{itemize}[nosep,leftmargin=*]
\item \textbf{Non-uniform allocation rarely helps.} At moderate compression (4$\times$), only 1/5 architectures benefits meaningfully. Four architectures show no improvement or minor degradation, with downstream evaluation confirming this pattern (Section~\ref{sec:cross_arch}).
\item \textbf{Unconstrained solvers systematically fail.} On \{2,4,8,16\}, both greedy and MCKP solvers over-allocate 2-bit heads due to compounding distortions, producing worse results than uniform 4-bit on all tested architectures (Section~\ref{sec:compounding}).
\item \textbf{A simple diagnostic predicts outcomes.} A two-metric check on head sensitivity distributions correctly classifies all tested models, telling practitioners whether to invest in profiling (Section~\ref{sec:diagnostic}).
\end{itemize}

We frame this work as a systematic empirical study---a resource contribution providing benchmarks, taxonomy, and practical guidelines.
We report all negative results transparently: downstream degradation on non-Spread architectures, inference latency overhead, K/V separation failures, and scalability limitations.

%% ========== RELATED WORK ==========
\section{Related Work}

\paragraph{KV cache quantization.}
KIVI~\cite{liu2024kivi} applies uniform per-channel quantization to keys and per-token quantization to values.
KVQuant~\cite{hooper2024kvquant} introduces sensitivity-weighted non-uniform quantization with outlier-aware calibration.
GEAR~\cite{kang2024gear} uses low-rank approximation combined with quantization.
These methods generally assume that more sophisticated allocation improves quality, but do not systematically evaluate when this assumption fails.

\paragraph{Rate--distortion allocation.}
RateQuant~\cite{zhang2024ratequant} formulates per-head allocation as a rate--distortion optimization, providing closed-form solutions for continuous bit-widths and greedy allocation for consecutive integers \{2,\ldots,6\}.
Their Proposition~6 guarantees greedy optimality on consecutive integer sets via polymatroid structure.
However, hardware-aligned bit-widths \{2,4,8,16\} are non-consecutive, breaking this guarantee.
RateQuant's Appendix~M acknowledges compounding effects but does not address them.

\paragraph{Mixed-precision quantization.}
HAWQ~\cite{dong2020hawq} and related methods use Hessian-based sensitivity for weight quantization precision selection.
Our work applies a similar profiling-then-allocating paradigm to KV cache quantization, but focuses on characterizing when this paradigm succeeds versus fails.

%% ========== METHODOLOGY ==========
\section{Methodology}\label{sec:method}

\subsection{Per-Head Sensitivity Profiling}

We profile each KV head independently by quantizing it to each candidate bit-width $b \in \{2,4,8,16\}$ while keeping all other heads at 16-bit, measuring the resulting perplexity delta $\Delta_h(b)$ on a calibration set.
Our high-resolution configuration uses 64 sequences of 2048 tokens from WikiText-2.

We verify that the sensitivity-based allocator automatically assigns higher bit-widths to the most sensitive heads without manual exclusion rules.
On Qwen2.5-7B, removing the Layer-0 exclusion constraint yields identical allocations (PPL 10.91 in both cases), confirming the profiler correctly identifies sensitive heads.
Cross-dataset calibration (WikiText-2 vs.\ C4) produces 96.4\% head-level agreement (108/112 heads), with a PPL difference of only 0.034---confirming that sensitivity profiles reflect model properties rather than calibration data artifacts.

\subsection{Constrained Allocation via MCKP}

Given per-head sensitivity profiles, we formulate bit-width allocation as a Multiple-Choice Knapsack Problem (MCKP):
\begin{equation}
\min_{\{b_h\}} \sum_h \Delta_h(b_h) \quad \text{s.t.} \quad \sum_h b_h \leq B, \quad |\{h : b_h = 2\}| \leq k
\end{equation}
where $B$ is the total bit budget and $k$ is a cardinality constraint limiting the number of 2-bit heads.
The cardinality constraint addresses a compounding failure mode detailed in Section~\ref{sec:compounding}.

\subsection{Sensitivity Distribution Taxonomy}\label{sec:taxonomy}

We classify architectures by their head sensitivity distribution shape:
\begin{itemize}[nosep,leftmargin=*]
\item \textbf{Spread}: Multiple heads with large cross-bitwidth gradients (e.g., Qwen2.5-7B: Layer-0 cluster of 4 heads with 26$\times$ dominance). Non-uniform allocation can exploit this structure.
\item \textbf{Spike}: Single extreme outlier dominates (e.g., Yi-1.5-9B: L4H3 holds 64\% of total delta). The solver assigns 8-bit to that head but has no other heads to differentiate.
\item \textbf{Flat}: Near-uniform sensitivity across heads (e.g., Mistral-7B: profiler reports near-zero differences). Non-uniform allocation produces zero signal.
\end{itemize}

%% ========== EXPERIMENTS ==========
\section{Experiments}\label{sec:experiments}

All experiments use WikiText-2 test set for perplexity evaluation unless otherwise noted. We use a unified BF16 evaluation pipeline with consistent parameters across all models.

\subsection{Cross-Architecture Benchmark}\label{sec:cross_arch}

Table~\ref{tab:cross_arch} presents the central result: per-head allocation outcomes across five architectures at 4$\times$ compression (budget = 0.25).

\begin{table}[h]
\centering
\caption{Cross-architecture per-head allocation benchmark (4$\times$ compression). ``Constrained vs.\ Uniform'' shows the relative PPL improvement of constrained MCKP over uniform 4-bit. Negative values indicate MCKP \emph{hurts}.}
\label{tab:cross_arch}
\small
\begin{tabular}{lcccccc}
\toprule
Model & GQA & KV Heads & Top-5\% Conc. & BF16 & Uniform 4-bit & MCKP ($k$) \\
\midrule
Qwen2.5-7B & 7:1 & 112 & 0.61 & 9.81 & 15.45 & \textbf{10.91} ($k\!=\!8$) \\
Yi-1.5-9B & 8:1 & 192 & 0.71 & --- & 7.22 & 7.22 ($k\!=\!8$) \\
Llama-3.1-8B & 4:1 & 256 & 0.57 & 8.09 & 8.47 & 8.45 ($k\!=\!2$) \\
Mistral-7B-v0.3 & 4:1 & 256 & 0.22 & 7.04 & 7.13 & 7.17 ($k\!=\!4$) \\
Gemma-2-9B & 2:1 & 336 & 0.23 & --- & 9.03 & 9.05 \\
\midrule
\multicolumn{4}{l}{\textit{Constrained vs.\ Uniform ($\Delta$)}} & & & \\
\multicolumn{4}{l}{Qwen: $+29.4\%$; Yi: $\approx 0\%$; Llama: $+0.2\%$; Mistral: $-0.4\%$; Gemma-2: $-0.3\%$} \\
\bottomrule
\end{tabular}
\end{table}

Only Qwen2.5-7B, classified as a Spread architecture, shows substantial benefit from non-uniform allocation.
Llama-3.1-8B shows a marginal +0.2\% improvement (0.02 PPL), which falls within calibration noise (multi-seed std = 0.065 on Qwen).
Mistral-7B and Gemma-2-9B both show that MCKP \emph{hurts}: the profiler reports near-zero sensitivity differences (concentration 0.22--0.23), producing allocations worse than uniform.
Yi-1.5-9B, despite having the highest GQA ratio (8:1) and highest concentration (0.71), shows zero benefit---its single-outlier Spike pattern cannot be exploited by MCKP.

\paragraph{Mistral k-ablation.}
On Mistral, PPL monotonically increases with $k$ (Table~\ref{tab:mistral_k}). Every additional 2-bit head hurts quality. The optimal allocation is $k\!=\!0$, i.e., uniform 4-bit.

\begin{table}[h]
\centering
\caption{Mistral-7B-v0.3 $k$-ablation (budget = 0.25). PPL increases monotonically---no value of $k$ improves over uniform.}
\label{tab:mistral_k}
\small
\begin{tabular}{lcccc}
\toprule
Config & PPL & 2-bit heads & 4-bit heads & Note \\
\midrule
$k=0$ (uniform) & 7.131 & 0 & 256 & \textbf{optimal} \\
$k=4$ & 7.165 & 4 & 252 & best constrained \\
$k=8$ & 7.200 & 8 & 247 & \\
$k=16$ & 7.311 & 16 & 238 & \\
$k=18$ & 7.323 & 18 & 235 & worst \\
\bottomrule
\end{tabular}
\end{table}

\paragraph{Llama allocation.}
On Llama-3.1-8B, MCKP naturally selects only 2 heads for 2-bit. The cardinality constraint ($k\!=\!2$) never binds.
The resulting allocation \{2:2, 4:253, 8:1\} yields PPL 8.45 vs.\ uniform 8.47---a 0.02 PPL difference that is not statistically distinguishable from noise without multi-seed confidence intervals on this architecture.

\subsection{Downstream Task Evaluation}\label{sec:downstream}

We evaluate on four downstream tasks (ARC-Challenge, HellaSwag, PIQA, WinoGrande) at 0-shot across five architectures (Table~\ref{tab:downstream}).

\begin{table}[h]
\centering
\caption{Downstream evaluation (4$\times$ compression, 0-shot). $\Delta$ is MCKP minus uniform 4-bit in percentage points.}
\label{tab:downstream}
\small
\begin{tabular}{lccccc}
\toprule
Model (Shape) & ARC-C $\Delta$ & HellaSwag $\Delta$ & PIQA $\Delta$ & WinoGrande $\Delta$ & Avg $\Delta$ \\
\midrule
Qwen (Spread) & +3.16pp & +4.00pp & +1.25pp & +1.97pp & \textbf{+1.9pp} \\
Yi (Spike) & $-$3.24pp & $-$0.60pp & $-$0.33pp & +1.27pp & $-$0.77pp \\
Llama (Spike) & $-$2.65pp & $-$0.14pp & $-$0.65pp & $-$0.15pp & $-$0.99pp \\
Mistral (Flat) & \multicolumn{4}{c}{(per-task breakdown omitted)} & $-$0.48pp \\
Gemma-2 (Flat) & \multicolumn{4}{c}{(per-task breakdown omitted)} & +0.14pp \\
\bottomrule
\end{tabular}
\end{table}

Downstream results confirm the perplexity findings.
Qwen shows consistent gains across all four tasks (+1.25 to +4.00pp), with 72.8\% average recovery of the BF16--uniform gap (BF16 ARC-C: 0.529, uniform: 0.481, MCKP: 0.514).
The four non-Spread architectures show minor systematic degradation (Yi $-$0.77pp, Llama $-$0.99pp, Mistral $-$0.48pp, Gemma-2 +0.14pp noise).
Notably, Llama shows degradation on all 5 tested tasks---the pattern is systematic, not task-specific noise.

ARC-Challenge shows the largest per-task sensitivity to allocation strategy across models (Yi $-$3.24pp, Llama $-$2.65pp), consistent with its reliance on multi-step reasoning where subtle KV cache errors accumulate.

\textbf{Limitation}: We do not provide multi-seed confidence intervals for downstream metrics; the per-task deltas (especially ARC-C $-$3.24pp on Yi) may be questioned for statistical significance.

\subsection{Budget Sweep}\label{sec:budget}

Table~\ref{tab:budget} shows Qwen2.5-7B performance across compression ratios.

\begin{table}[h]
\centering
\caption{Qwen2.5-7B budget sweep: MCKP ($k\!=\!8$) vs.\ uniform at three compression ratios.}
\label{tab:budget}
\small
\begin{tabular}{lccc}
\toprule
Compression & MCKP PPL & Uniform PPL & Relative Improvement \\
\midrule
2$\times$ (budget 0.50) & 9.85 & 10.2 & +3.4\% \\
4$\times$ (budget 0.25) & 10.91 & 15.45 & +29.4\% \\
8$\times$ (budget 0.125) & 57.6 & 2644 & $>$45$\times$ \\
\bottomrule
\end{tabular}
\end{table}

At mild compression (2$\times$), the benefit is marginal (+3.4\%).
At moderate compression (4$\times$), non-uniform allocation provides substantial improvement---but only on Qwen, the single Spread architecture.
At aggressive compression (8$\times$), both methods produce high perplexity, but the gap is dramatic: MCKP prevents the catastrophic collapse that uniform 2-bit suffers.
The aggressive regime is arguably trivially expected---any non-uniform scheme will outperform uniform 2-bit when many heads are highly sensitive.

\subsection{Supporting Analysis: Compounding Effects}\label{sec:compounding}

We observe that unconstrained per-head allocation systematically fails on \{2,4,8,16\}.
Table~\ref{tab:compounding} shows a sweep of the cardinality parameter $k$ on Qwen2.5-7B, revealing a compounding pattern where perplexity degrades super-linearly as more heads are assigned 2-bit.

\begin{table}[h]
\centering
\caption{Qwen2.5-7B compounding sweep (budget = 0.25). Ratio is PPL relative to $k\!=\!0$.}
\label{tab:compounding}
\small
\begin{tabular}{lccl}
\toprule
$k$ (max 2-bit heads) & PPL & Ratio & Note \\
\midrule
0 & 9.81 & 1.00 & all $\geq$4-bit \\
8 & 10.05 & 1.02 & constraint threshold \\
16 & 10.42 & 1.06 & \\
32 & 11.60 & 1.18 & onset of rapid degradation \\
64 & 23.22 & 2.37 & catastrophic \\
96 & $\sim$65 & $\sim$6.6 & \\
112 (unconstrained) & $\sim$98.75 & $\sim$10.07 & \\
\bottomrule
\end{tabular}
\end{table}

The compounding pattern is architecture-specific.
Qwen exhibits a U-shaped PPL curve with optimal $k\!=\!8$.
Llama shows a flat-then-cliff pattern with the optimal being uniform ($k\!=\!0$) or very low $k$.
Mistral shows monotonic increase---every 2-bit head hurts.
This heterogeneity across architectures undermines any claim of a universal compounding threshold; instead, it supports the taxonomy in Section~\ref{sec:taxonomy}, where different distribution shapes predict different compounding behaviors.

We note that compounding of quantization errors is a known phenomenon---Zhang et al.~\cite{zhang2024ratequant} discuss it in their Appendix~M.
Our contribution is the systematic quantification across architectures, not the discovery of the phenomenon itself.

\subsection{Solver Comparison}\label{sec:solver}

Table~\ref{tab:solver} compares RateQuant's greedy algorithm (adapted to \{2,4,8,16\}) with our MCKP formulation, both unconstrained and constrained.

\begin{table}[h]
\centering
\caption{Solver comparison on \{2,4,8,16\} bit-widths (Qwen2.5-7B, budget = 0.25).}
\label{tab:solver}
\small
\begin{tabular}{lcc}
\toprule
Method & PPL & Bit Distribution \\
\midrule
Greedy unconstrained & 32.45 & \{2:68, 4:30, 8:4, 16:10\} \\
MCKP unconstrained & 36.92 & --- \\
\midrule
Greedy + $k\!=\!8$ & 11.23 & \{2:8, 4:100, 8:4\} \\
MCKP + $k\!=\!8$ & 11.23 & \{2:8, 4:100, 8:4\} \\
\midrule
Uniform 4-bit & 15.45 & all 4-bit \\
\bottomrule
\end{tabular}
\end{table}

Both unconstrained solvers fail catastrophically (PPL 32--37 vs.\ uniform 15.45), confirming that the problem lies in unconstrained allocation on non-consecutive bit-widths, not in solver design.
Under the cardinality constraint, both solvers produce \textbf{identical} allocations and PPL ($\Delta = 0.001$).
This demonstrates that the constraint reduces the feasible set to where solver choice is irrelevant---the constrained allocation problem is well-behaved.

We note that RateQuant's greedy is designed for consecutive integers \{2,\ldots,6\} where Proposition~6 guarantees optimality.
Evaluating it on \{2,4,8,16\} violates its stated assumptions.
The comparison is not a criticism of RateQuant but rather evidence that hardware-aligned non-consecutive bit-widths introduce difficulties that their framework does not address.

\subsection{MCKP vs.\ Oracle}

To quantify solver quality, we compare MCKP allocations against brute-force oracle on a subset of 8 layers (Table~\ref{tab:oracle}).

\begin{table}[h]
\centering
\caption{MCKP vs.\ oracle (Qwen2.5-7B, $k\!=\!8$, 8-layer subset).}
\label{tab:oracle}
\small
\begin{tabular}{lcc}
\toprule
Metric & Value & Note \\
\midrule
MCKP PPL & 10.91 & --- \\
Best oracle gap & +0.018 PPL & single 2-bit head swap \\
Average oracle gap & +0.005 PPL & --- \\
Layers where MCKP = Oracle & 5/8 & --- \\
Total oracle evaluations & 54 & --- \\
\bottomrule
\end{tabular}
\end{table}

The near-zero oracle gap (avg 0.005 PPL) confirms that MCKP is near-optimal under the constraint.
This also implies that the constrained problem is small enough that any reasonable solver suffices---consistent with the greedy--MCKP equivalence in Table~\ref{tab:solver}.

\subsection{L0-Heuristic Baseline}\label{sec:l0_heuristic}

We evaluate the simplest possible non-uniform strategy: assign 8-bit to all Layer-0 heads, uniform 4-bit to everything else.
This zero-cost heuristic requires no profiling whatsoever.

\begin{table}[h]
\centering
\caption{L0 heuristic vs.\ MCKP across architectures.}
\label{tab:l0_heuristic}
\small
\begin{tabular}{lcccc}
\toprule
Model & MCKP PPL & L0 Heuristic (median) & Uniform 4-bit & Profiling Cost \\
\midrule
Qwen2.5-7B & 10.91 & 11.52 & 15.45 & 65 min (MCKP) / 0 (L0) \\
Llama-3.1-8B & 8.45 & 8.58 & 8.47 & 173 min (MCKP) / 0 (L0) \\
Mistral-7B-v0.3 & 7.16 & --- & 7.13 & MCKP worse than uniform \\
\bottomrule
\end{tabular}
\end{table}

On Qwen, the L0 heuristic captures approximately 80\% of the MCKP benefit (PPL 11.52 vs.\ 10.91) at zero profiling cost.
The remaining 0.61 PPL gap represents the marginal value of full profiling and constrained optimization.
On Llama, the L0 heuristic (8.58) is worse than both MCKP (8.45) and uniform (8.47), reinforcing that on Spike/Flat architectures, even this simple heuristic adds complexity without benefit.
This result suggests that for practitioners operating on Spread-type models, protecting Layer-0 heads is the most cost-effective first step.

\subsection{Profiling Resolution}\label{sec:profiling}

We compare low-resolution (16 sequences $\times$ 512 tokens) and high-resolution (64 $\times$ 2048) profiling on Qwen2.5-7B.

\begin{table}[h]
\centering
\caption{Profiling resolution comparison (Qwen2.5-7B, $k\!=\!8$).}
\label{tab:profiling}
\small
\begin{tabular}{lccc}
\toprule
Resolution & MCKP PPL & 2-bit Candidates & Time \\
\midrule
Low-res (16$\times$512) & 11.23 & 37 & $\sim$5 min \\
High-res (64$\times$2048) & 10.91 & 100 & 65.5 min \\
Uniform 4-bit & 15.45 & --- & 0 min \\
\bottomrule
\end{tabular}
\end{table}

High-resolution profiling improves PPL by 0.32 (11.23 $\to$ 10.91) but changes only 8 of 112 head assignments.
The low-resolution profile already identifies the correct constraint ($k\!=\!8$).
The marginal improvement from high-resolution is real but modest relative to the 13$\times$ increase in profiling time, suggesting that low-resolution profiling is sufficient for the diagnostic decision (Section~\ref{sec:diagnostic}) while high-resolution is warranted only when full optimization is justified.

\subsection{K/V Separation Ablation}

We test whether allocating bit-widths independently for key and value heads improves over joint allocation.

\begin{table}[h]
\centering
\caption{K/V separation ablation (Qwen2.5-7B).}
\label{tab:kv_sep}
\small
\begin{tabular}{lcc}
\toprule
Method & PPL & Bit Distribution \\
\midrule
Joint K+V MCKP & 10.91 & \{2:8, 4:100, 8:4\} \\
Separate K/V MCKP & 11.22 & K:\{2:16,\ldots\}, V: all-4bit \\
Uniform 4-bit & 15.45 & all 4-bit \\
\bottomrule
\end{tabular}
\end{table}

K/V separation \emph{hurts} by 0.31 PPL.
Under separation, the solver assigns all value heads to uniform 4-bit and concentrates 2-bit assignments on key heads (16 at 2-bit), suggesting that value quantization sensitivity is relatively uniform.
This negative result indicates that our framework does not naturally handle K/V asymmetry---an area where RateQuant's approach (asymmetric K/V budgets) may be more appropriate.

\subsection{Multi-Seed Calibration Stability}

We evaluate allocation stability across 3 calibration seeds on Qwen2.5-7B (Table~\ref{tab:multiseed}).

\begin{table}[h]
\centering
\caption{Multi-seed calibration (Qwen2.5-7B, $k\!=\!8$).}
\label{tab:multiseed}
\small
\begin{tabular}{lcc}
\toprule
Seed & PPL & Note \\
\midrule
Seed 0 & 11.02 & \\
Seed 1 & 10.90 & \\
Seed 2 & 11.05 & \\
\midrule
Mean $\pm$ std & 10.99 $\pm$ 0.065 & \\
3-way allocation agreement & 90.2\% & \\
\bottomrule
\end{tabular}
\end{table}

The PPL variance is small (std = 0.065), and the 95\% CI [10.86, 11.12] remains well below uniform (15.45), confirming that the Qwen result is robust.
However, 3-way allocation agreement of 90.2\% means approximately 10\% of heads receive different bit-width assignments depending on the calibration seed.
This instability---while not impacting PPL significantly---raises concerns for deployment reproducibility.

\textbf{Limitations}: (1) Three seeds is the minimum; standard practice recommends 5--10 seeds with 95\% CI.
(2) We do not provide multi-seed results for other architectures; the marginal results on Llama ($\Delta = +0.2\%$) and Mistral ($\Delta = -0.4\%$) cannot be confirmed as statistically distinguishable from zero without seed-level confidence intervals on those models.

\subsection{Cross-Dataset Robustness}

\begin{table}[h]
\centering
\caption{Cross-dataset calibration (Qwen2.5-7B).}
\label{tab:cross_dataset}
\small
\begin{tabular}{lccc}
\toprule
Calibration $\to$ Eval & MCKP PPL & Uniform PPL & Improvement \\
\midrule
WT2 $\to$ WT2 & 10.91 & 15.45 & +29.4\% \\
C4 $\to$ WT2 & 11.29 & 15.45 & +27.0\% \\
WT2 $\to$ C4 & 20.56 & 35.42 & +41.9\% \\
\midrule
Allocation agreement (WT2 vs.\ C4) & \multicolumn{3}{c}{96.4\% (108/112 heads)} \\
\bottomrule
\end{tabular}
\end{table}

Cross-dataset calibration produces near-identical allocations (96.4\% agreement).
The C4-calibrated allocation evaluated on WikiText-2 achieves +27.0\% improvement, close to the matched-calibration +29.4\%.
This confirms that head sensitivity is a stable model property rather than a calibration data artifact.

\subsection{Inference Performance}\label{sec:runtime}

\begin{table}[h]
\centering
\caption{Inference performance (Qwen2.5-7B, single A100).}
\label{tab:runtime}
\small
\begin{tabular}{lcccc}
\toprule
Config & Throughput (tok/s) & Decode Latency (ms) & KV Memory & vs.\ Uniform 4-bit \\
\midrule
BF16 & 29.6 & 33.78 & 21 MB & --- \\
Uniform 4-bit & 23.6 & 42.46 & 5.25 MB & --- \\
MCKP ($k\!=\!8$) & 21.8 & 46.09 & 5.25 MB & $-$7.6\% / +8.5\% \\
\bottomrule
\end{tabular}
\end{table}

Mixed-precision MCKP achieves the same 4$\times$ memory compression as uniform 4-bit but incurs an additional \textbf{8.5\% decode latency overhead} (46.09 vs.\ 42.46 ms) and \textbf{7.6\% throughput reduction} (21.8 vs.\ 23.6 tok/s).
This overhead stems from heterogeneous bit-width dequantization kernels that are not optimized for mixed-precision access patterns.
For 4/5 architectures where MCKP provides no PPL benefit, this latency overhead is a \emph{pure loss}.
Even on Qwen, practitioners must weigh the +29.4\% PPL improvement against the 8\% throughput penalty---a tradeoff that depends on deployment constraints.
Kernel optimization to reduce mixed-precision overhead is an important direction for future work.

%% ========== DIAGNOSTIC ==========
\section{When to Profile: A Practitioner Diagnostic}\label{sec:diagnostic}

Based on the cross-architecture analysis, we propose a simple two-metric diagnostic for deciding whether per-head profiling is worthwhile:
\begin{enumerate}[nosep]
\item \textbf{Top-5\% concentration} $\gtrsim 0.50$: the top 5\% of heads account for at least half of total sensitivity.
\item \textbf{Top-1/Top-5 ratio} $t_1/t_5 \lesssim 0.70$: sensitivity is distributed across multiple top heads (Spread), not concentrated in a single outlier (Spike).
\end{enumerate}

Both conditions must hold to recommend non-uniform allocation. Table~\ref{tab:diagnostic} shows the diagnostic's predictions.

\begin{table}[h]
\centering
\caption{Diagnostic classification of tested architectures.}
\label{tab:diagnostic}
\small
\begin{tabular}{lccccl}
\toprule
Model & Conc. & $t_1/t_5$ & Shape & Diagnostic & Actual Outcome \\
\midrule
Qwen2.5-7B & 0.61 & 0.57 & Spread & Profile & \textbf{+29.4\%} \\
Yi-1.5-9B & 0.71 & 0.90 & Spike & Skip & $\approx$0\% \\
Llama-3.1-8B & 0.57 & 0.85 & Spike & Skip & +0.2\% (noise) \\
Mistral-7B-v0.3 & 0.22 & --- & Flat & Skip & $-$0.4\% \\
Gemma-2-9B & 0.23 & 0.38 & Flat & Skip & $-$0.3\% \\
\bottomrule
\end{tabular}
\end{table}

The diagnostic correctly classifies all 5 tested architectures: it recommends profiling only for Qwen (the sole beneficiary) and correctly skips all others.
Yi-1.5 is the key illustration: despite high concentration (0.71), its high $t_1/t_5$ ratio (0.90) reveals a Spike pattern where a single outlier dominates, making per-head optimization futile.

\textbf{Limitations}: The diagnostic is validated on $n\!=\!5$ architectures.
While the pattern is consistent, $n\!=\!5$ provides directional evidence, not a statistical claim of generalization.
The thresholds (0.50, 0.70) are derived from this sample and may require adjustment as more architectures are tested.

%% ========== SCALABILITY ==========
\section{Limitations}\label{sec:limitations}

\paragraph{Scalability.}
Our 70B-scale profiling attempt (Qwen2.5-72B) was killed due to GPU OOM/timeout.
All reported results are on 7--14B models.
Extrapolating profiling cost from 65 minutes for 7B/112 heads to 72B/1024+ heads suggests 10+ hours of GPU time, which is impractical for deployment workflows.
Efficient profiling strategies (e.g., layer-wise subsampling, proxy-model transfer) are essential future work.

\paragraph{Architecture coverage.}
Five architectures provide a consistent pattern but do not constitute statistical proof.
The diagnostic thresholds are empirically derived and may not transfer to all model families.
Additional architectures (e.g., Falcon, DeepSeek, InternLM) are needed to validate generality.

\paragraph{Inference overhead.}
The 8\% latency overhead from mixed-precision dequantization is not inherent to the method but reflects unoptimized kernel implementations.
Without dedicated mixed-precision kernels, the practical value of non-uniform allocation is diminished even when PPL improves.

\paragraph{Downstream evaluation.}
We evaluate on four 0-shot tasks without multi-seed confidence intervals for downstream metrics.
Per-task deltas on non-Spread architectures (e.g., Yi ARC-C $-$3.24pp) warrant further investigation with proper statistical testing.

\paragraph{K/V asymmetry.}
Our framework treats K and V heads as joint items.
The K/V separation ablation shows this hurts (0.31 PPL degradation), suggesting that specialized K/V handling---as in RateQuant---may be more appropriate.

%% ========== CONCLUSION ==========
\section{Conclusion}

We present a systematic empirical study of per-head KV cache quantization on hardware-aligned bit-widths.
Our key finding is sobering: at moderate compression, non-uniform allocation benefits only 1 of 5 tested architectures.
For the remaining four, uniform 4-bit is as good or better, and the profiling overhead and mixed-precision latency penalty represent pure cost.

We contribute a sensitivity-based taxonomy (Spread, Spike, Flat) and a two-metric diagnostic that correctly predicts allocation outcomes on all tested models, enabling practitioners to skip unnecessary profiling.
We additionally characterize a compounding failure mode in unconstrained allocation and show that a cardinality constraint prevents it, reducing the problem to where solver choice is irrelevant.

We report all negative results transparently: downstream degradation on non-Spread architectures, 8\% inference latency overhead, K/V separation failures, profiling scalability limitations, and 10\% allocation instability across calibration seeds.
These results challenge the implicit assumption that non-uniform allocation always improves KV cache quantization and provide practitioners with the data needed to make informed deployment decisions.

\bibliographystyle{plain}
\begin{thebibliography}{9}

\bibitem{hooper2024kvquant}
C.~Hooper, S.~Kim, H.~Mohammadzadeh, M.~W.~Mahoney, Y.~S.~Shao, K.~Keutzer, and A.~Gholami.
\newblock KVQuant: Towards 10 million context length LLM inference with KV cache quantization.
\newblock \emph{arXiv preprint arXiv:2401.18079}, 2024.

\bibitem{liu2024kivi}
Z.~Liu, J.~Yuan, H.~Jin, S.~Zhong, Z.~Xu, V.~Braverman, B.~Chen, and X.~Hu.
\newblock KIVI: A tuning-free asymmetric 2bit quantization for KV cache.
\newblock \emph{arXiv preprint arXiv:2402.02750}, 2024.

\bibitem{zhang2024ratequant}
Z.~Zhang, X.~Xiao, and Z.~Lin.
\newblock RateQuant: A rate-distortion framework for KV cache quantization.
\newblock \emph{arXiv preprint}, 2024.

\bibitem{kang2024gear}
H.~Kang, Q.~Zhang, S.~Kundu, G.~Jeong, Z.~Liu, T.~Krishna, and T.~Zhao.
\newblock GEAR: An efficient KV cache compression recipe for near-lossless generative inference of LLM.
\newblock \emph{arXiv preprint arXiv:2403.05527}, 2024.

\bibitem{dong2020hawq}
Z.~Dong, Z.~Yao, A.~Gholami, M.~W.~Mahoney, and K.~Keutzer.
\newblock HAWQ: Hessian AWare Quantization of neural networks with mixed-precision.
\newblock In \emph{ICCV}, 2019.

\end{thebibliography}

\end{document}
